%% KAIRIS working manuscript for Electrochimica Acta.
%% Based on Elsevier CAS single-column template v2.4 (LPPL).
\documentclass[a4paper,fleqn]{cas-sc}

\usepackage[numbers,sort&compress]{natbib}
\usepackage{booktabs}
\usepackage{siunitx}

\begin{document}
\let\WriteBookmarks\relax
\def\floatpagepagefraction{1}
\def\textpagefraction{.001}

\shorttitle{KAIRIS: a reproducible multiperspective EIS protocol}
\shortauthors{Author list pending}

\title[mode=title]{KAIRIS (Kramers--Kronig-Assessed Impedance, Relaxation,
and Integrated Spectroscopy): A Multiperspective Protocol for Reproducible
Electrochemical Impedance Spectroscopy}

\author[1]{Author details pending}
\cormark[1]
\ead{corresponding.author@example.org}
\credit{Conceptualization, Methodology, Software, Validation, Writing}

\affiliation[1]{organization={Affiliation pending},
  city={City pending},
  country={Country pending}}

\cortext[1]{Corresponding author}

\begin{abstract}
Electrochemical impedance spectroscopy can provide complementary views of
interfacial, transport and distributed processes, but its interpretation is
limited by data admissibility, non-unique representations and workflows that do
not preserve uncertainty or disagreement. This study introduces KAIRIS
(Kramers--Kronig-Assessed Impedance, Relaxation, and Integrated Spectroscopy), a
reproducible protocol integrating spectral validation, distribution of
relaxation times, global spectral features and electrochemical spectral
cartography. The confirmatory design, quantitative results and externally
validated domain of applicability will be inserted only after the frozen
analysis is completed.
\end{abstract}

\begin{highlights}
\item KAIRIS integrates Lin--KK, DRT, GSF and ESC in a reproducible EIS workflow.
\item Distinct spectral observables are compared without forcing artificial agreement.
\item Deterministic rules preserve abstention without requiring a language model.
\end{highlights}

\begin{keywords}
electrochemical impedance spectroscopy \sep distribution of relaxation times
\sep Kramers--Kronig \sep reproducibility \sep spectral analysis
\end{keywords}

\maketitle

\section{Introduction}\label{sec:introduction}

% Define the EIS interpretation problem, the limits of equivalent-circuit
% uniqueness, and the need for complementary observables and abstention.

\section{Scientific questions and hypotheses}\label{sec:hypotheses}

% State the frozen H1--H6 hypotheses and distinguish retrospective from
% prospective validation.

\section{Methods}\label{sec:methods}

\subsection{KAIRIS protocol}\label{sec:protocol}

\subsection{Spectral admissibility and Lin--KK}\label{sec:linkk}

\subsection{Distribution of relaxation times}\label{sec:drt}

\subsection{Global spectral features}\label{sec:gsf}

\subsection{Electrochemical spectral cartography}\label{sec:esc}

\subsection{Datasets and leakage prevention}\label{sec:datasets}

\subsection{External 316L confirmatory cohort}\label{sec:external}

The selected external cohort contains 42 impedance spectra of welded 316L
stainless steel measured in aerated \SI{0.1}{\mole\per\litre} NaCl. It comprises
eight surface/region conditions and two passivity reference groups. Selection,
exclusions, endpoints and the analysis lock will be described from the registered
external-validation addendum before results are accessed.

\subsection{Blinded expert assessment}\label{sec:experts}

\subsection{Baselines, ablations and sensitivity}\label{sec:baselines}

\subsection{Statistical analysis}\label{sec:statistics}

\subsection{Software and computational reproducibility}\label{sec:software}

\section{Results}\label{sec:results}

% Do not write confirmatory results before the analysis lock is opened.

\subsection{Engineering integrity and portability}

\subsection{Synthetic parameter recovery}

\subsection{External scientific validation}

\subsection{Expert agreement and abstention}

\subsection{Ablations, sensitivity and discordance taxonomy}

\section{Discussion}\label{sec:discussion}

\subsection{Complementarity of DRT, GSF and ESC}

\subsection{Domain of applicability and limitations}

\section{Conclusions}\label{sec:conclusions}

\section*{Data and code availability}

Code is developed openly at
\url{https://github.com/Emanuel-963/KAIRIS}. Persistent release and data
identifiers will be inserted before submission.

\section*{CRediT authorship contribution statement}

Author roles pending confirmation.

\section*{Declaration of competing interest}

Declaration pending author confirmation.

\section*{Acknowledgements}

Funding and acknowledgements pending confirmation.

\printcredits

\bibliographystyle{cas-model2-names}
\bibliography{references}

\end{document}
